\section{Conclusion}
\label{sec:conclusion}

In shallow full-space Penalty-X QAOA, poor feasible sampling did not have a
single cause. Exact depth nesting exposed run-specific optimizer inadequacy,
while matched lower-energy/lower-feasibility states showed that optimizer
repair does not resolve objective misalignment. O2 then demonstrated that the
same ansatz retained feasible-concentration capacity not induced by
mean-energy training.

This makes the O0$\rightarrow$O2$\rightarrow$O3 sequence the central result.
On graph-disjoint held-out data, CVaR-0.10 improved feasibility compensation
over O0 and met the preregistered non-inferiority criterion against the
exact-feasibility O2 capacity control. CVaR therefore used the available
energy distribution more effectively for feasible entry in this controlled
setting; O2 itself remains mechanistic and non-deployable.

The robustness evidence sets equally important boundaries. Greater depth can
improve terminal feasibility while increasing certified optimization
failures. Fixed-endpoint sampling is more stable than finite-shot retraining,
whose intervals cross zero. The $m=20$ reversal and administrative $m=22$
censoring rule out extrapolative scaling claims.

These synthetic RCSP tasks are classically tractable, and the study tests
neither structured mixers, objective-specific initialization, hardware noise,
real networks, nor quantum advantage. Its practical message is diagnostic:
constrained-QAOA evaluations should distinguish representation dilution,
ansatz capacity, optimizer adequacy, objective alignment, and sampling effects
before treating a poor endpoint as a fundamental algorithmic limitation.
